\documentclass[aspectratio=169]{beamer}
\usetheme{metropolis}

\usepackage[utf8]{inputenc}
\usepackage[spanish]{babel}
\usepackage{graphicx}
\usepackage{xcolor}
\usepackage{tikz}
\usetikzlibrary{positioning, arrows.meta, babel}
\usepackage{tcolorbox}
\tcbuselibrary{skins, breakable}
\usepackage{fontawesome5}
\usepackage{amsmath}

% ─────────────────────────────────────────────────────────────
%  PALETA
% ─────────────────────────────────────────────────────────────
\definecolor{mainColor}{HTML}{00a499}
\definecolor{accentColor}{HTML}{EA7600}
\definecolor{bgLight}{HTML}{E6F7F6}
\definecolor{codeGray}{HTML}{2B2B2B}

% ─────────────────────────────────────────────────────────────
%  METROPOLIS
% ─────────────────────────────────────────────────────────────
\setbeamercolor{normal text}{fg=codeGray, bg=white}
\setbeamercolor{frametitle}{bg=mainColor, fg=white}
\setbeamercolor{alerted text}{fg=accentColor}
\setbeamercolor{progress bar}{fg=accentColor, bg=mainColor!25}
\setbeamercolor{block title}{bg=mainColor, fg=white}
\setbeamercolor{block body}{bg=bgLight, fg=codeGray}
\setbeamercolor{block title alerted}{bg=accentColor, fg=white}
\setbeamercolor{block body alerted}{bg=accentColor!10, fg=codeGray}
\metroset{progressbar=frametitle, sectionpage=none, numbering=fraction}
\usepackage{hyperref}

% ─────────────────────────────────────────────────────────────
%  CAJA DE CONCEPTO
% ─────────────────────────────────────────────────────────────
\newtcolorbox{conceptbox}[1]{
  colback=bgLight, colframe=mainColor,
  title={\small #1}, fonttitle=\bfseries,
  coltitle=white, colbacktitle=mainColor,
  arc=4pt, boxrule=0.6pt,
  left=4pt, right=4pt, top=2pt, bottom=2pt
}

\newcommand{\tech}[1]{\texttt{\textcolor{mainColor}{#1}}}

\hypersetup{colorlinks=true, urlcolor=mainColor}

% ─────────────────────────────────────────────────────────────
%  PORTADA
% ─────────────────────────────────────────────────────────────
\title{Modelos de Inteligencia Artificial en Negocios\\{\normalsize |}\\ Inteligencia Artificial Aplicada a la Industria\\{\tiny Clase 3: Datos para Entrenamiento}}
\subtitle{{\small Semestre 1, 2026 · Universidad de Santiago de Chile}}
\author{\href{https://www.linkedin.com/in/byron-caices/}{Byron Caices Lima} \\ \small Departamento de Ingeniería Industrial}
\date{\today}

\begin{document}

\maketitle

% =============================================================
%  AGENDA
% =============================================================

\begin{frame}{¿Qué Veremos Hoy?}
  \small
  \begin{columns}[T]
    \begin{column}{0.32\textwidth}
      \begin{conceptbox}{\faBrain\enspace Parte 1}
        \small
        \textbf{¿Qué es un modelo de IA?}\\[4pt]
        La idea de \textit{aprender con ejemplos} y por qué los datos son su materia prima.
      \end{conceptbox}
    \end{column}
    \begin{column}{0.32\textwidth}
      \begin{conceptbox}{\faTable\enspace Parte 2}
        \small
        \textbf{Tipos de datos}\\[4pt]
        Tabulares vs. no estructurados. Por qué la IA siempre los \textbf{traduce a números}.
      \end{conceptbox}
    \end{column}
    \begin{column}{0.32\textwidth}
      \begin{conceptbox}{\faFlask\enspace Parte 3}
        \small
        \textbf{Laboratorio}\\[4pt]
        Preparar nuestro CSV limpio para que una futura IA pueda aprender de él.
      \end{conceptbox}
    \end{column}
  \end{columns}
  \vspace{0.3cm}
  \begin{alertblock}{\faLightbulb\enspace Objetivo de la sesión}
    \centering\small
    Entender que \textbf{datos limpios todavía no son suficientes} --- y aprender a transformarlos para que la IA pueda usarlos.
  \end{alertblock}
\end{frame}

% =============================================================
%  PUENTE CLASE 2
% =============================================================

\begin{frame}{¿Dónde Quedamos en la Clase 2?}
  \begin{columns}[T]
    \begin{column}{0.55\textwidth}
      \small
      Recibimos un CSV \textbf{sucio} de defectos de manufactura y lo limpiamos:
      \vspace{0.15cm}
      \begin{itemize}\setlength\itemsep{4pt}
        \item Convertimos \tech{defect\_date} de texto a fecha real.
        \item Rellenamos los nulos en \tech{severity} e \tech{inspection\_method}.
        \item Eliminamos las filas con \tech{repair\_cost} negativo.
      \end{itemize}

      \vspace{0.3cm}
      \textbf{Pregunta:} ¿podemos usar ese CSV directamente para que una IA aprenda de él?
    \end{column}
    \begin{column}{0.42\textwidth}
      \begin{alertblock}{\faExclamationTriangle\enspace Aún no}
        \small
        Los datos limpios todavía contienen:
        \begin{itemize}\setlength\itemsep{2pt}
          \item Categorías como ``Structural'' o ``Visual'' (\textbf{texto, no números})
          \item Columnas con escalas muy distintas
          \item Un solo bloque de datos sin separar
        \end{itemize}
      \end{alertblock}
      \vspace{0.15cm}
      \begin{conceptbox}{\faArrowRight\enspace Hoy}
        \small\centering
        Aprendemos a \textbf{transformar} los datos para que la IA pueda procesarlos.
      \end{conceptbox}
    \end{column}
  \end{columns}
\end{frame}

% =============================================================
%  PARTE 1: ¿QUÉ ES IA / ML?
% =============================================================

% ─────────────────────────────────────────────────────────────
%  ¿QUÉ ES ML?
% ─────────────────────────────────────────────────────────────
\begin{frame}{Antes de Empezar: ¿Qué es Machine Learning?}
  \begin{columns}[T]
    \begin{column}{0.55\textwidth}
      \small
      Cuando hablamos de IA aplicada a la industria, en la mayoría de los casos nos referimos a \textbf{Machine Learning (ML)} --- ``aprendizaje automático''.

      \vspace{0.2cm}
      \begin{conceptbox}{\faLightbulb\enspace La idea central}
        \small
        Un programa tradicional sigue \textbf{reglas escritas a mano}.\\[4pt]
        Un modelo de Machine Learning \textbf{descubre las reglas por sí mismo} a partir de muchos ejemplos.
      \end{conceptbox}
    \end{column}
    \begin{column}{0.42\textwidth}
      \begin{block}{\faCode\enspace Programa tradicional}
        \tiny\ttfamily
        si costo $>$ 500:\\
        \quad marcar como ``caro''\\
        si\_no:\\
        \quad marcar como ``barato''
      \end{block}
      \vspace{0.1cm}
      \begin{alertblock}{\faRobot\enspace Modelo de ML}
        \small
        ``Aquí tienes 1000 defectos pasados con su costo real. \textbf{Aprende tú} qué características hacen que un defecto sea caro.''
      \end{alertblock}
    \end{column}
  \end{columns}
  \vspace{0.2cm}
  \begin{conceptbox}{\faBookOpen\enspace La analogía clave}
    \centering\small
    Los datos son a la IA \textbf{lo que los libros de texto son a un estudiante}.\\
    Si el material está desordenado o lleno de errores, el estudiante no aprende bien.
  \end{conceptbox}
\end{frame}

% ─────────────────────────────────────────────────────────────
%  EJEMPLO HUMANO -> MODELO
% ─────────────────────────────────────────────────────────────
\begin{frame}{Aprender con Ejemplos: La Idea de ``Entrenar''}
  \begin{columns}[T]
    \begin{column}{0.5\textwidth}
      \small
      Cuando un humano aprende a reconocer una manzana, no lee un manual: ve \textbf{muchas manzanas} y poco a poco capta los patrones.

      \vspace{0.2cm}
      Un modelo de IA hace algo similar: necesita \textbf{muchos ejemplos} del fenómeno que queremos que aprenda.

      \vspace{0.2cm}
      Cada ejemplo tiene dos partes:
      \begin{itemize}\setlength\itemsep{3pt}
        \item \textbf{Características}: la información que describe el caso.
        \item \textbf{Resultado}: lo que pasó realmente.
      \end{itemize}
    \end{column}
    \begin{column}{0.46\textwidth}
      \begin{conceptbox}{\faIndustry\enspace En nuestro caso}
        \small
        \textbf{Características de cada defecto:}\\[2pt]
        \tech{defect\_type}, \tech{defect\_location}, \tech{severity}, \tech{inspection\_method}\\[8pt]
        \textbf{Resultado correspondiente:}\\[2pt]
        \tech{repair\_cost} (cuánto costó realmente repararlo)
      \end{conceptbox}
      \vspace{0.15cm}
      \begin{alertblock}{\faLightbulb\enspace La meta}
        \small
        Si le mostramos miles de ejemplos, el modelo aprende: \textit{``cuando un defecto es Estructural y crítico, el costo suele ser alto''}.
      \end{alertblock}
    \end{column}
  \end{columns}
\end{frame}

% ─────────────────────────────────────────────────────────────
%  VOCABULARIO X / y
% ─────────────────────────────────────────────────────────────
\begin{frame}{Vocabulario que Vamos a Usar}
  \small
  \begin{columns}[T]
    \begin{column}{0.48\textwidth}
      \begin{block}{\faList\enspace Las características --- ``X''}
        \small
        \begin{itemize}\setlength\itemsep{3pt}
          \item Las columnas que el modelo usa como \textbf{pistas}.
          \item Se escribe en \textbf{mayúscula} porque suele ser una matriz (varias columnas).
          \item También se llaman \textbf{features} (en inglés).
        \end{itemize}
      \end{block}
    \end{column}
    \begin{column}{0.48\textwidth}
      \begin{alertblock}{\faBullseye\enspace El objetivo --- ``y''}
        \small
        \begin{itemize}\setlength\itemsep{3pt}
          \item La columna que el modelo intenta \textbf{adivinar}.
          \item Se escribe en \textbf{minúscula} porque es una sola columna.
          \item También se llama \textbf{target} (en inglés).
        \end{itemize}
      \end{alertblock}
    \end{column}
  \end{columns}
  \vspace{0.25cm}
  \begin{conceptbox}{\faArrowRight\enspace La meta del modelo}
    \centering\small
    Aprender la relación: \textbf{X $\rightarrow$ y}\\[3pt]
    \textit{``Dadas estas características, ¿cuál sería el resultado?''}
  \end{conceptbox}
\end{frame}

% =============================================================
%  PARTE 2: TIPOS DE DATOS
% =============================================================

% ─────────────────────────────────────────────────────────────
%  TABULARES VS NO ESTRUCTURADOS
% ─────────────────────────────────────────────────────────────
\begin{frame}{Datos Tabulares vs. Datos No Estructurados}
  \small
  \begin{columns}[T]
    \begin{column}{0.48\textwidth}
      \begin{block}{\faTable\enspace Datos Tabulares}
        \small
        Filas y columnas, como una tabla de Excel.
        \vspace{4pt}
        \begin{itemize}\setlength\itemsep{2pt}
          \item Cada columna tiene un significado claro.
          \item Formato: CSV, Excel, base SQL.
        \end{itemize}
        \vspace{4pt}
        \textbf{Ejemplos en industria:}
        \begin{itemize}\setlength\itemsep{1pt}
          \item Registros de defectos \checkmark
          \item Ventas mensuales
          \item Lecturas de sensores
        \end{itemize}
      \end{block}
    \end{column}
    \begin{column}{0.48\textwidth}
      \begin{block}{\faImages\enspace Datos No Estructurados}
        \small
        No caben en una tabla. Cada ``unidad'' es un archivo complejo.
        \vspace{4pt}
        \begin{itemize}\setlength\itemsep{2pt}
          \item Difíciles de procesar directamente.
          \item Formato: imagen, audio, video, texto libre.
        \end{itemize}
        \vspace{4pt}
        \textbf{Ejemplos en industria:}
        \begin{itemize}\setlength\itemsep{1pt}
          \item Fotos de productos defectuosos
          \item Audio de máquinas
          \item Reclamos de clientes
        \end{itemize}
      \end{block}
    \end{column}
  \end{columns}
  \vspace{0.2cm}
  \begin{alertblock}{\faLightbulb\enspace La regla universal de la IA}
    \centering\small
    Cualquiera sea el tipo de dato, el modelo lo recibe \textbf{como una matriz de números}.\\
    Las imágenes, el texto y las tablas: \textbf{todo se traduce a números} antes de entrar al modelo.
  \end{alertblock}
\end{frame}

% =============================================================
%  PARTE 3: TRANSFORMACIONES
% =============================================================

% ─────────────────────────────────────────────────────────────
%  PROBLEMA 1: TEXTO → NÚMEROS
% ─────────────────────────────────────────────────────────────
\begin{frame}{Problema 1: Las Máquinas No Entienden Palabras}
  \small
  La primera fila de nuestro dataset se ve más o menos así:

  \vspace{0.2cm}
  \centering
  {\footnotesize
  \begin{tabular}{|l|l|l|l|}
    \hline
    \textbf{defect\_type} & \textbf{defect\_location} & \textbf{severity} & \textbf{inspection\_method} \\\hline
    Structural & Surface & Critical & Visual \\\hline
  \end{tabular}}

  \vspace{0.3cm}
  \begin{columns}[T]
    \begin{column}{0.48\textwidth}
      \begin{alertblock}{\faExclamationTriangle\enspace El problema}
        \small
        Un modelo de IA solo trabaja con \textbf{números}. La palabra ``Critical'' no significa nada para él.
      \end{alertblock}
    \end{column}
    \begin{column}{0.48\textwidth}
      \begin{conceptbox}{\faCogs\enspace La solución}
        \small
        Traducir cada categoría a un número, sin inventar un orden falso entre ellas.\\[3pt]
        Técnica estándar: \textbf{One-Hot Encoding}.
      \end{conceptbox}
    \end{column}
  \end{columns}
\end{frame}

% ─────────────────────────────────────────────────────────────
%  ONE-HOT ENCODING EXPLICADO
% ─────────────────────────────────────────────────────────────
\begin{frame}{One-Hot Encoding: Una Columna por Categoría}
  \small
  La idea es simple: por cada categoría posible creamos una nueva columna que vale \textbf{1 si la fila pertenece a esa categoría, 0 si no}.

  \vspace{0.25cm}
  \centering
  {\footnotesize
  \begin{tabular}{|c||c|c|c|}
    \hline
    \textbf{defect\_type} & \textbf{Structural} & \textbf{Functional} & \textbf{Cosmetic} \\\hline\hline
    Structural & \textbf{1} & 0 & 0 \\\hline
    Functional & 0 & \textbf{1} & 0 \\\hline
    Cosmetic & 0 & 0 & \textbf{1} \\\hline
  \end{tabular}}

  \vspace{0.3cm}
  \begin{columns}[T]
    \begin{column}{0.48\textwidth}
      \begin{alertblock}{\faTimesCircle\enspace ¿Por qué no asignar 1, 2, 3?}
        \small
        Si decimos \textit{Structural=1, Functional=2, Cosmetic=3}, el modelo cree que \textbf{Cosmetic $>$ Functional $>$ Structural}.\\[3pt]
        Eso es un \textbf{orden falso}: son categorías sin ranking.
      \end{alertblock}
    \end{column}
    \begin{column}{0.48\textwidth}
      \begin{conceptbox}{\faPython\enspace En Python: 1 línea}
        \small
        \tech{pd.get\_dummies(X)}\\[4pt]
        Pandas detecta las columnas de texto y aplica One-Hot Encoding automáticamente.
      \end{conceptbox}
    \end{column}
  \end{columns}
\end{frame}

% ─────────────────────────────────────────────────────────────
%  EJEMPLO PASO A PASO: ONE-HOT ENCODING
% ─────────────────────────────────────────────────────────────
\begin{frame}{Ejemplo Paso a Paso: One-Hot Encoding en Nuestro Dataset}
  \small
  \textbf{Partamos con 3 defectos reales de nuestro dataset:}
  \vspace{0.2cm}

  \begin{columns}[T]
    \begin{column}{0.48\textwidth}
      \textbf{ANTES (con categorías de texto):}
      \vspace{4pt}
      \par{\footnotesize
      \begin{tabular}{|l|l|}
        \hline
        \textbf{defect\_type} & \textbf{severity} \\\hline
        Structural & Minor \\
        Cosmetic & Critical \\
        Functional & Moderate \\\hline
      \end{tabular}}

      \vspace{0.3cm}
      \small Las máquinas no entienden ``Structural'' ni ``Critical''.
    \end{column}
    \begin{column}{0.48\textwidth}
      \textbf{DESPUÉS (solo números 0 y 1):}
      \vspace{4pt}
      \par{\tiny
      \begin{tabular}{|c|c|c|c|c|c|}
        \hline
        \textbf{Struct} & \textbf{Funct} & \textbf{Cosm} & \textbf{Minor} & \textbf{Crit} & \textbf{Mod} \\\hline
        1 & 0 & 0 & 1 & 0 & 0 \\
        0 & 0 & 1 & 0 & 1 & 0 \\
        0 & 1 & 0 & 0 & 0 & 1 \\\hline
      \end{tabular}}

      \vspace{0.15cm}
      \small Las máquinas entienden perfectamente los 0s y 1s.
    \end{column}
  \end{columns}

  \vspace{0.3cm}
  \begin{alertblock}{\faLightbulb\enspace La lógica del 0 y 1}
    \centering\small
    Para cada fila: \textbf{exactamente una columna vale 1}, el resto valen 0.\\
    ¿Por qué? Porque cada defecto \textbf{es una sola cosa a la vez}.\\
    No puede ser Structural Y Cosmetic simultáneamente.
  \end{alertblock}
\end{frame}

% ─────────────────────────────────────────────────────────────
%  ¿POR QUÉ NO ASIGNAR 1, 2, 3 DIRECTAMENTE?
% ─────────────────────────────────────────────────────────────
\begin{frame}{¿Por Qué No Simplemente Asignar 1, 2, 3?}
  \small
  \textbf{Propuesta tentadora:} asignar Structural=1, Functional=2, Cosmetic=3.
  \vspace{0.2cm}

  \begin{columns}[T]
    \begin{column}{0.48\textwidth}
      \begin{alertblock}{\faTimesCircle\enspace ¿Cuál es el problema?}
        \small
        El modelo interpretaría estos números como una \textbf{escala real}:\\[4pt]
        ``Cosmetic (3) es MÁS que Functional (2),\\
        y Functional (2) es MÁS que Structural (1).''\\[4pt]
        \textbf{Pero eso es falso.} Son solo categorías distintas, sin jerarquía.
      \end{alertblock}
    \end{column}
    \begin{column}{0.48\textwidth}
      \begin{block}{\faCheckCircle\enspace La solución: One-Hot}
        \small
        Con One-Hot, cada categoría es \textbf{independiente}:
        \vspace{4pt}
        \par{\tiny
        \begin{tabular}{|c|c|c|}
          \hline
          \textbf{S} & \textbf{F} & \textbf{C} \\\hline
          1 & 0 & 0 \\
          0 & 1 & 0 \\
          0 & 0 & 1 \\\hline
        \end{tabular}}

        \vspace{2pt}
        \par No hay orden. El modelo no puede confundir ``3 es mayor que 1''.
      \end{block}
    \end{column}
  \end{columns}

  \vspace{0.2cm}
  \begin{conceptbox}{\faKey\enspace Regla general}
    \centering\small
    \textbf{Siempre} que tengas categorías sin un orden natural → usa One-Hot Encoding.
  \end{conceptbox}
\end{frame}

% ─────────────────────────────────────────────────────────────
%  PROBLEMA 2: ESCALADO
% ─────────────────────────────────────────────────────────────
\begin{frame}{Problema 2: Cuando los Números Viven en Escalas Distintas}
  \begin{columns}[T]
    \begin{column}{0.5\textwidth}
      \small
      Imaginen un dataset con dos columnas:
      \vspace{0.15cm}
      \begin{itemize}\setlength\itemsep{2pt}
        \item \tech{edad\_maquina} $\rightarrow$ entre 1 y 30 (años)
        \item \tech{horas\_uso\_total} $\rightarrow$ entre 0 y 50.000 (horas)
      \end{itemize}

      \vspace{0.25cm}
      Para muchos modelos, la columna con \textbf{números más grandes pesa más} en sus decisiones --- aunque las dos variables sean igualmente importantes.

      \vspace{0.2cm}
      \textit{Es como medir personas en metros y árboles en kilómetros: el árbol parece insignificante \textbf{solo por la unidad}.}
    \end{column}
    \begin{column}{0.46\textwidth}
      \begin{conceptbox}{\faRulerCombined\enspace La solución: escalar}
        \small
        Llevar todas las columnas numéricas a un \textbf{rango común}, por ejemplo entre 0 y 1.
        \vspace{4pt}
        \par\textbf{Min-Max Scaling:}
        \[ x' = \frac{x - \min}{\max - \min} \]
        Después, todos los valores quedan entre 0 y 1, sin importar la escala original.
      \end{conceptbox}
      \vspace{0.1cm}
      \begin{alertblock}{\faPython\enspace En sklearn}
        \small\centering
        \tech{MinMaxScaler()}
      \end{alertblock}
    \end{column}
  \end{columns}
\end{frame}

% ─────────────────────────────────────────────────────────────
%  PROBLEMA 3: ENTRENAR vs EVALUAR
% ─────────────────────────────────────────────────────────────
\begin{frame}{Problema 3: ¿Cómo Saber si el Modelo Realmente Aprendió?}
  \begin{columns}[T]
    \begin{column}{0.55\textwidth}
      \small
      Imaginen que un profesor enseña matemáticas con \textbf{100 ejercicios resueltos}, y luego evalúa al alumno con \textbf{esos mismos 100 ejercicios}.

      \vspace{0.2cm}
      ¿Aprendió de verdad o solo \textbf{memorizó las respuestas}?

      \vspace{0.3cm}
      Lo mismo pasa con los modelos de IA:
      \vspace{0.1cm}
      \begin{itemize}\setlength\itemsep{3pt}
        \item Si los entrenamos con \textbf{todos los datos}\ldots
        \item \ldots y luego los evaluamos con \textbf{esos mismos datos}\ldots
        \item \ldots no sabremos si el modelo aprendió patrones reales.
      \end{itemize}
    \end{column}
    \begin{column}{0.42\textwidth}
      \begin{conceptbox}{\faCut\enspace La solución}
        \small
        Separar los datos \textbf{antes de entrenar} en dos grupos:
        \vspace{4pt}
        \begin{itemize}\setlength\itemsep{2pt}
          \item \textbf{Entrenamiento (~80\%)}: el modelo estudia con ellos.
          \item \textbf{Prueba (~20\%)}: nunca los ve durante el aprendizaje. Sirven para evaluarlo después.
        \end{itemize}
      \end{conceptbox}
      \vspace{0.1cm}
      \begin{alertblock}{\faPython\enspace En sklearn}
        \small\centering
        \tech{train\_test\_split(X, y, test\_size=0.2)}
      \end{alertblock}
    \end{column}
  \end{columns}
\end{frame}

% ─────────────────────────────────────────────────────────────
%  VISUAL DEL SPLIT
% ─────────────────────────────────────────────────────────────
\begin{frame}{Visualización del Split}
  \centering
  \vspace{0.3cm}

  \begin{tikzpicture}
    \fill[mainColor!75] (0,0) rectangle (10,1.2);
    \fill[accentColor!75] (10,0) rectangle (12.5,1.2);
    \draw[thick] (0,0) rectangle (12.5,1.2);
    \node[white, font=\large\bfseries] at (5,0.6) {ENTRENAMIENTO --- 80\%};
    \node[white, font=\large\bfseries] at (11.25,0.6) {PRUEBA --- 20\%};

    \node[font=\small] at (5,-0.4) {El modelo aprende aquí};
    \node[font=\small] at (11.25,-0.4) {Evaluamos al final};
  \end{tikzpicture}

  \vspace{0.8cm}
  \begin{columns}[T]
    \begin{column}{0.48\textwidth}
      \begin{conceptbox}{\faGraduationCap\enspace Datos de entrenamiento}
        \small
        Son los \textbf{ejemplos resueltos} que le mostramos al modelo. Aquí el modelo descubre patrones.
      \end{conceptbox}
    \end{column}
    \begin{column}{0.48\textwidth}
      \begin{alertblock}{\faClipboardCheck\enspace Datos de prueba}
        \small
        Son el \textbf{examen}. El modelo nunca los vio antes. Si responde bien acá, podemos confiar en él.
      \end{alertblock}
    \end{column}
  \end{columns}
\end{frame}

% =============================================================
%  PARTE 4: CASO HOY
% =============================================================

% ─────────────────────────────────────────────────────────────
%  NUESTRO CASO HOY
% ─────────────────────────────────────────────────────────────
\begin{frame}{Nuestro Caso de Hoy}
  \begin{columns}[T]
    \begin{column}{0.55\textwidth}
      \small
      \textbf{Situación:} Seguimos en la planta industrial. Los datos limpios de la Clase 2 nos dejan en condiciones de plantear una idea:

      \vspace{0.15cm}
      \textit{``¿Y si más adelante construimos una IA que estime el costo de reparación de un defecto futuro a partir de sus características?''}

      \vspace{0.2cm}
      \textbf{Hoy NO entrenaremos el modelo todavía.} Hoy preparamos los datos para que esa próxima etapa sea posible.

      \vspace{0.2cm}
      \textbf{Variables del problema:}
      \begin{itemize}\setlength\itemsep{2pt}
        \item Características (X): tipo, ubicación, severidad, método de inspección
        \item Objetivo (y): \tech{repair\_cost}
      \end{itemize}
    \end{column}
    \begin{column}{0.42\textwidth}
      \begin{conceptbox}{\faFileExport\enspace Entregable de la clase}
        \small
        Cuatro archivos CSV listos para entrenamiento:
        \vspace{4pt}
        \begin{itemize}\setlength\itemsep{2pt}
          \item \tech{X\_train.csv}
          \item \tech{y\_train.csv}
          \item \tech{X\_test.csv}
          \item \tech{y\_test.csv}
        \end{itemize}
      \end{conceptbox}
      \vspace{0.15cm}
      \begin{alertblock}{\faIndustry\enspace ¿Por qué este paso es crítico?}
        \small\centering
        Sin esta preparación, los algoritmos de la próxima unidad \textbf{simplemente no funcionarían}.
      \end{alertblock}
    \end{column}
  \end{columns}
\end{frame}

% ─────────────────────────────────────────────────────────────
%  HOJA DE RUTA DEL LABORATORIO
% ─────────────────────────────────────────────────────────────
\begin{frame}{Hoy en el Laboratorio: Hoja de Ruta}
  \centering
  \vspace{0.15cm}

  \begin{tikzpicture}[
      node distance=0.30cm and 0.18cm,
      step/.style={rectangle, rounded corners=4pt, draw=mainColor, fill=bgLight,
                   text width=2.0cm, align=center, font=\tiny\bfseries, minimum height=1.1cm},
      arr/.style={->, thick, color=accentColor, shorten >=2pt, shorten <=2pt}
    ]
    \node[step] (s1) {\faDownload\\[3pt]1. Cargar\\datos limpios};
    \node[step, right=of s1] (s2) {\faBullseye\\[3pt]2. Definir\\X e y};
    \node[step, right=of s2] (s3) {\faCode\\[3pt]3. One-Hot\\Encoding};
    \node[step, right=of s3] (s4) {\faRulerCombined\\[3pt]4. Escalar\\valores};
    \node[step, right=of s4] (s5) {\faCut\\[3pt]5. Dividir\\train / test};
    \node[step, right=of s5] (s6) {\faSave\\[3pt]6. Exportar\\4 CSV};

    \draw[arr] (s1) -- (s2);
    \draw[arr] (s2) -- (s3);
    \draw[arr] (s3) -- (s4);
    \draw[arr] (s4) -- (s5);
    \draw[arr] (s5) -- (s6);
  \end{tikzpicture}

  \vspace{0.3cm}
  \begin{columns}[T]
    \begin{column}{0.48\textwidth}
      \begin{conceptbox}{\faPython\enspace Herramientas}
        \small
        \begin{itemize}\setlength\itemsep{3pt}
          \item \tech{pandas}, \tech{numpy}, \tech{matplotlib}
          \item \tech{scikit-learn} \textbf{(nuevo)}
        \end{itemize}
      \end{conceptbox}
    \end{column}
    \begin{column}{0.48\textwidth}
      \begin{conceptbox}{\faKey\enspace Conceptos clave}
        \small
        \begin{itemize}\setlength\itemsep{2pt}
          \item Características (X) y objetivo (y)
          \item One-Hot Encoding
          \item Escalado de valores
          \item Train / Test Split
        \end{itemize}
      \end{conceptbox}
    \end{column}
  \end{columns}

  \vspace{0.2cm}
  \begin{alertblock}{\faArrowRight\enspace Próxima unidad}
    \centering\small
    Con estos archivos preparados podremos ver cómo \textbf{los modelos de Machine Learning aprenden} de los datos.
  \end{alertblock}

  \vspace{0.2cm}
  \centering
  \large\textcolor{mainColor}{\faArrowRight}\enspace
  \textbf{\normalsize ¡Abran Google Colab!}
\end{frame}

\end{document}
